\section{Validación y Verificación (Objetivo Específico 4)}

\subsection{Entornos Análogos de Prueba}
Las pruebas se realizarán en entornos que simulen la morfología lunar:
\begin{itemize}
    \item \textbf{Cantera de arena/grava:} Para validar tracción y hundimiento.
    \item \textbf{Circuito de obstáculos:} Para validar la detección y evasión autónoma.
    \item \textbf{Zona de pendientes:} Para validar límites de estabilidad y par motor.
\end{itemize}

\subsection{Plan de Pruebas de Sistema}
\begin{table}[H]
    \centering
    \begin{tabularx}{\textwidth}{|l|l|X|X|}
        \hline
        \textbf{ID Prueba} & \textbf{Req. Asociado} & \textbf{Procedimiento} & \textbf{Criterio de Éxito} \\ \hline
        TP-01 & RF-NAV-01 & El rover debe avanzar 5m con 3 obstáculos aleatorios. & No colisiona y llega al destino. \\ \hline
        TP-02 & RNF-PER-01 & Medir tiempo de recorrido en una recta de 10m. & Velocidad media $\ge$ 5 cm/s. \\ \hline
        TP-03 & RF-ENE-02 & Forzar descarga de batería simulada. & El rover activa modo seguro a 15\%. \\ \hline
    \end{tabularx}
\end{table}

\subsection{Matriz de Trazabilidad}
Esta matriz asegura que cada objetivo del proyecto esté cubierto por requerimientos, arquitectura y pruebas.

\begin{table}[H]
    \centering
    \begin{tabularx}{\textwidth}{|X|l|l|l|}
        \hline
        \textbf{Objetivo Específico} & \textbf{Requerimiento} & \textbf{Subsistema} & \textbf{ID Prueba} \\ \hline
        1. Definir requerimientos & RF-NAV-01 & Navegación & TP-01 \\ \hline
        2. Diseñar arquitectura & RF-ENE-01 & Potencia & TP-03 \\ \hline
        3. Construir prototipo & RNF-PER-01 & Movilidad & TP-02 \\ \hline
    \end{tabularx}
\end{table}
